% =============================================================================
\chapter{Eksperymenty wstępne i wnioski}
\label{ch:preliminary}
% =============================================================================

Benchmark i system ewaluacyjny opisane w kolejnych rozdziałach nie były naszym pierwszym projektem.
Dwa kierunki pilotażowe zakończyły się niepowodzeniem, a przyczyny tych niepowodzeń ukształtowały
ostateczną metodykę. Opisujemy je, ponieważ są to przydatne wyniki negatywne dla każdego, kto planuje
podobne badanie.

\section{Zadania w stylu HumanEval są zbyt łatwe}
\label{sec:humaneval}

\paragraph{Motywacja} Naturalnym punktem wyjścia był HumanEval~\cite{chen2021codex}. Jest to
najczęściej raportowany benchmark generowania kodu, jego problemy są z założenia niezależne od
języka, a MultiPL-E~\cite{cassano2023multiple} pokazał, jak przetłumaczyć je wraz z testami na inne
języki. Planowaliśmy przenieść podzbiór jego problemów do siedmiu języków badania i porównać odsetki
rozwiązań.

\paragraph{Obserwacja} Problemy okazały się zdecydowanie zbyt łatwe. W naszym pilotażu modele
rozwiązały praktycznie każdy problem za pierwszym podejściem w każdym języku, w tym w językach
o małych zasobach. Wyniki wszystkich języków znalazły się na pułapie lub tuż pod nim, więc nie mogły
wykazać żadnych różnic między paradygmatami. Nie dało się też badać samonaprawy ani
idiomatyczności, ponieważ nie było czego naprawiać, a rozwiązania były zbyt krótkie, by ujawnić
jakikolwiek projekt.
\todo{dodać liczby z pilotażu: modele, liczba problemów i języków, odsetki rozwiązań}

\paragraph{Przyczyny} Z perspektywy czasu efekt sufitowy był do przewidzenia, z czterech powodów.
\begin{enumerate}
  \item \emph{Rozmiar problemu.} Typowy problem HumanEval wymaga jednej krótkiej funkcji na
    listach, napisach lub liczbach całkowitych. Taka funkcja ma kilka linii w dowolnym języku i nie
    wymaga ani struktur danych, ani zarządzania własnością pamięci, więc nie zostawia miejsca na
    dekompozycję specyficzną dla paradygmatu.
  \item \emph{Nasycenie.} Obecne czołowe modele rozwiązują zdecydowaną większość problemów
    HumanEval, dlatego społeczność przeszła do trudniejszych lub nowszych
    benchmarków~\cite{liu2023evalplus,jain2024livecodebench,jimenez2024swebench}.
  \item \emph{Kontaminacja.} Problemy i wiele rozwiązań są publiczne od 2021 roku i prawdopodobnie
    znajdują się w danych treningowych obecnych modeli~\cite{riddell2024contamination}. Wysoki wynik
    może wtedy świadczyć o zapamiętaniu, a nie o umiejętności programowania.
  \item \emph{Słabe testy.} Przy niewielkiej liczbie testów na problem przechodzą rozwiązania
    częściowo błędne~\cite{liu2023evalplus}. Przez to pułap wydaje się jeszcze wyższy, niż jest
    w rzeczywistości.
\end{enumerate}

\paragraph{Konsekwencje dla projektu} Postanowiliśmy napisać własne zadania, spełniające poniższe
wymagania, z których każde odpowiada na jeden z powyższych problemów:
\begin{description}[leftmargin=2.4em,labelwidth=2em,itemsep=2pt]
  \item[D1] \emph{Kontrakty z wieloma funkcjami.} Każde zadanie wymaga dwóch powiązanych algorytmów
    za niewielkim API z własnymi typami danych, np.\ grafem, drzewem AST lub B-drzewem, tak aby
    rozwiązanie trzeba było zaprojektować, a nie tylko napisać.
  \item[D2] \emph{Precyzyjne specyfikacje.} Reguły rozstrzygania remisów, wartości błędów
    i przypadki brzegowe są wyspecyfikowane tak, aby oczekiwany wynik był jednoznaczny i dało się
    go porównać dokładnie (na przykład dokładny kształt B-drzewa po ciągu wstawień i usunięć).
  \item[D3] \emph{Mocne testy.} Każde zadanie ma 20 testów (10 na algorytm), w tym dane puste
    i zdegenerowane, niepoprawne argumenty oraz testy obciążeniowe na dużych danych.
  \item[D4] \emph{Ukryte testy.} Model nigdy nie widzi kodu testów, więc nie może dopasować
    rozwiązania do testów.
  \item[D5] \emph{Bezpieczeństwo pamięci.} Tam, gdzie język na to pozwala, testy są uruchamiane pod
    Valgrindem, więc wycieki i niepoprawne dostępy do pamięci liczą się jako błędy.
  \item[D6] \emph{Kompozycje trudniejsze do zapamiętania.} Klasyczne algorytmy są łączone
    z nietypowymi wymaganiami, takimi jak ciągłe wykrywanie kolizji oparte na sumie Minkowskiego,
    selekcja BFPRT raportująca ślad wyboru pivotów czy mediana dwóch posortowanych ciągów przy
    ścisłym limicie odczytów, aby zmniejszyć wartość zapamiętanych rozwiązań.
\end{description}
\Cref{ch:benchmark} opisuje powstały benchmark. \Cref{ch:results} pokazuje, że unika on efektu
sufitowego: za pierwszym podejściem rozwiązano tylko 80\,\% zadań, a błędy różnią się między
językami.

\section{Lokalne modele o otwartych wagach są niepraktyczne}
\label{sec:local-llms}

\paragraph{Motywacja} Drugim planem było lokalne uruchamianie modeli o otwartych wagach. Obiecywało
to zerowy koszt krańcowy przebiegu, pełną kontrolę nad dekodowaniem (temperatura, ziarna losowe),
stałe wagi modelu zapewniające powtarzalność oraz niezależność od komercyjnego API.
\todo{podać nazwy modeli lokalnych (rodzina, rozmiar, kwantyzacja), środowisko inferencji i sprzęt}

\paragraph{Przeszkoda 1: wydajność} Generowanie na dostępnym nam sprzęcie było zbyt wolne jak na
badanie tej skali. Pojedyncza próba składa się z długiego promptu (specyfikacja wraz z kontraktem,
a w rundach naprawy także poprzedni kod i log testów) i długiej odpowiedzi (kompletny moduł
o długości 100--280 linii, zob.\ \cref{tab:efficiency}). Konwersacja naprawcza rośnie z każdą
rundą. Dla porównania: w naszym ostatecznym przebiegu \haiku{} generował średnio około 6\,500
tokenów wyjściowych na wywołanie, łącznie z rozumowaniem, i 635\,269 tokenów wyjściowych w całym
przebiegu. Przy przykładowej lokalnej szybkości dekodowania 10~tokenów/s zajęłoby to około
11~minut na wywołanie i prawie 18~godzin samego dekodowania na jedno przejście przez benchmark, bez
przetwarzania promptów i bez powtórzonych próbek, których wymaga badanie statystyczne. Mniejsze lub
silniej skwantyzowane modele, które działały szybciej, pisały wyraźnie gorszy kod, więc mierzyłyby
ograniczenia sprzętu, a nie wpływ paradygmatu.
\todo{dodać zmierzoną przepustowość (tokeny/s) i czas rzeczywisty na zadanie z pilotażu}

\paragraph{Przeszkoda 2: modele nie zwracały samego kodu} System zapisuje odpowiedź modelu
bezpośrednio do pliku źródłowego, kompiluje go i uruchamia testy. Wymaga to, aby odpowiedź była
kodem i niczym więcej. Modele lokalne często się do tego nie stosowały, nawet gdy prompt mówił o tym
wprost. Dodawały wyjaśnienia przed kodem i po nim, otaczały go znacznikami bloków Markdown, podawały
kilka alternatywnych wersji, zwracały fragmenty zamiast kompletnych modułów albo powtarzały części
promptu. Każda taka odpowiedź kończyła się błędem kompilacji, który mierzył wykonywanie instrukcji,
a nie umiejętność programowania. Co gorsza, efekt ten nie był jednakowy dla wszystkich języków,
więc mieszał się ze zmienną, którą chcieliśmy badać. Heurystyki wycinające kod z prozy (wybór
największego bloku kodu, usunięcie tekstu wstępnego) rozwiązywały proste przypadki. Zawodziły,
gdy bloków kodu było kilka lub brakowało znaczników, a każda ręczna decyzja, który blok jest
„odpowiedzią”, czyniłaby wyniki mniej obiektywnymi. Ostatecznie wyniki tych przebiegów były bardzo
trudne do interpretacji i zrezygnowaliśmy w tym badaniu z modeli lokalnych.

\paragraph{Decyzja} Przeszliśmy na hostowane modele Claude dostępne przez API Anthropic. Ten wybór
usuwa obie przeszkody. Przepustowość jest wystarczająca, by przeprowadzić cały benchmark w około
trzy kwadranse czasu pracy modelu (\cref{sec:cost}). Modele oferują też \emph{adaptacyjne
myślenie} (ang.~\emph{adaptive thinking})~\cite{anthropic2026docs}, które zwraca rozumowanie
w osobnych blokach \emph{thinking}, dzięki czemu odpowiedź tekstowa może być czystym kodem. Sędzia
korzysta z \emph{ustrukturyzowanych odpowiedzi} (ang.~\emph{structured outputs}), tj.\ schematu
JSON, który odpowiedź musi spełniać. Zachowaliśmy awaryjny mechanizm wycinania kodu i każdą
odpowiedź, która go wymagała, rejestrujemy jako \emph{naruszenie formatu}. Dzięki temu
wykonywanie instrukcji staje się mierzoną wielkością, a nie ukrytym źródłem szumu
(\cref{sec:extraction}).

\begin{findingbox}[Wniosek.]
Zgodność z wymaganym formatem odpowiedzi jest warunkiem wstępnym zautomatycznej ewaluacji kodu, a nie
szczegółem. Nawet model hostowany otoczył kod znacznikami bloków Markdown w 37,8\,\% odpowiedzi,
a w dwóch odpowiedziach wplótł do pliku źródłowego tekst niebędący kodem
(\cref{sec:cost,sec:failure-analysis}). Benchmark powinien zatem mierzyć i raportować zgodność
formatu oddzielnie od poprawności.
\end{findingbox}
